\section{Exercitii}

\subsection{Statistici fara bucla}
\label{problem:statistici-vectorizate}

\href{notebooks/statistici-vectorizate.ipynb}{notebook}
$\bullet$
\hyperref[code:statistici-vectorizate]{surse}

Calculam medie, deviatie standard si normalizare pe un vector, intai cu
bucla \pycode{for} ca sa vedem exact ce calculeaza, apoi vectorizat cu
NumPy, si comparam timpul de executie intre variante.

\subsection{Suma si maximul, cu NumPy}
\label{problem:suma-maxim-numpy}

\href{notebooks/suma-maxim-numpy.ipynb}{notebook}
$\bullet$
\hyperref[code:suma-maxim-numpy]{surse}

Acelasi \hyperref[problem:suma-vector]{suma} si
\hyperref[problem:maxim-vector]{maxim} din Curs 0, de data asta cu
\pycode{v.sum()} si \pycode{v.max()}. Compara codul cu bucla de mana pe
care ai scris-o acolo, calculul e identic, doar scris o singura data,
in libraria.

\subsection{Sortare si elemente unice}
\label{problem:sortare-unice}

\href{notebooks/sortare-unice.ipynb}{notebook}
$\bullet$
\hyperref[code:sortare-unice]{surse}

\pycode{np.sort()} sorteaza un vector, \pycode{np.unique()} scoate
duplicatele (si il sorteaza pe deasupra), \pycode{np.argmax()} iti da
pozitia maximului, nu valoarea lui. Trei functii pe care o sa le
folosesti des la prelucrarea datelor.

% TODO: extindem lista de exercitii pana la 10-20, cate un pic din
% fiecare directie. Idei, de adaugat pe rand, dupa ce confirmam ca
% formatul de mai sus e ok:
%
% NumPy: np.cumsum, np.where (selectie conditionata fara bucla),
%   broadcasting intre doua array-uri de dimensiuni diferite,
%   reshape / array 2D (paralela cu matricile din C++), slicing
%   multidimensional, masti booleene combinate cu & si |
%
% Pandas: citire si filtrare pe un CSV real (nu inventat), groupby cu
%   mai multe agregari deodata, sort_values, merge intre doua tabele,
%   valori lipsa (isna / fillna)
